\chapter{Continuous hardware, discrete silicon}
\label{ch:pillarii}

The structural theorem of this chapter is corpus-safe (\cid{19}) --- not a
finished lecture.
A plain-English restatement lives in \corpus{directives/theorem-1-plain.md};
a Swift-free check is \texttt{python3 Scripts/tensorlut\_math\_ref.py}, and the
smallest worked example of the LUT view is
\texttt{python3 Scripts/toy\_cipher\_demo.py} (\corpus{INTRO.md}).
The separable melt--freeze--snap certificate is corpus-safe (\cid{44}).
Shatter / hold grades and campaign empirics remain seminar-depth and are
\emph{not} the theorem~\cite{helut-tensorlut-thm}.

\growswhen{A completeness theorem for \emph{multi-LUT topological} melt:
that melt recovers cryptographic structure combinatoric sieves miss, for a
stated class of netlists beyond a fully observed 1-LUT.
Theorem~1$''$ (\cid{44}) is the separable interpolant, not that class.
Stream-cipher melts remain on the trajectory.}

\section{The problem boolean search cannot see}

A combinatoric sieve on a cipher netlist proposes discrete genotypes:
INIT bits, plugboard pairs, rotor orders.
On a short ciphertext the search space is both enormous and full of
linguistic hallucinations (the P1030680 campaign is the case study, not
the claim).
Boolean search \emph{rejects}.
It does not \emph{propose} a nearby reciprocal structure that is almost a
stecker.

Pillar~II melts the hardware~\cite{helut-paper,helut-tensorlut-thm}.

\section{Setup}

Let $L$ be the number of LUT6 cells and $w\in[0,1]^{64L}$ the concatenated
INIT genome.
Let $y(w)$ be the \emph{multilinear extension} of those INITs on $[0,1]$
(exact on $\{0,1\}$ inputs).
Let $t$ be a target bit vector of matching width.
Soft crypto fitness and discreteness penalty:
\begin{align}
  F_{\mathrm{crypto}}(w)
  &= -\lVert y(w)-t\rVert_2^2, \\
  \pi(w)
  &= \sum_i w_i(1-w_i), \\
  F(w)
  &= F_{\mathrm{crypto}}(w)-\lambda\,\pi(w),
  \qquad \lambda\ge 0.
\end{align}

\begin{remark}[Implementation schedule]
HELUT's GA uses a squeeze $\lambda=\lambda_{\max}(p')^2$.
Theorem~1 needs only $\lambda\ge 0$.
\end{remark}

Emitter: $E(w)_i=\mathbf{1}[w_i\ge\tfrac12]$.
Freeze mask $M\subseteq\{1,\dots,64L\}$ removes frozen coordinates from
$\pi$.
Stecker genotypes are \emph{partial involutions} (disjoint pairs) on
$\{0,\dots,25\}$.

Every Yosys \texttt{\$lut} of width $1\ldots 6$ pads to a 64-wide
\texttt{Float32} INIT.
The forward pass is branchless and GPU-shaped.
It is not torus FHE.

\section{Theorem 1 (structural)}
\label{sec:thm-tensorlut}

\reproduced{C19}{TensorLUT continuous$\to$discrete Theorem~1:
six machine-checked lemmas
(\texttt{TensorLUTFormal.certificate()}).
Not a cryptanalytic break.
Statement: \corpus{directives/tensorlut-theorem.md}.
Reproduce:
\texttt{swift test -c release --filter testTensorLUTFormalCertificate}.}

\begin{theorem}[TensorLUT continuous$\to$discrete; \corpusthm{1}]
\label{thm:tensorlut}
Assume the hypotheses on the certificate: the LUT is multilinear; the GA
mutates only unfrozen $w_i\in[0,1]$; the involution sandwich freezes core
INITs.
Then:
\begin{enumerate}
  \item \textbf{Discreteness.}
        $\pi(w)\ge 0$, with equality iff $w\in\{0,1\}^{64L}$.
  \item \textbf{Crypto MSE.}
        $F_{\mathrm{crypto}}(w)\le 0$, with equality iff $y(w)=t$.
  \item \textbf{Combined objective.}
        $F(w)\le F_{\mathrm{crypto}}(w)$; if $\pi(w)>0$, increasing
        $\lambda$ strictly decreases $F$.
  \item \textbf{Emitter.}
        $E$ is idempotent on $\{0,1\}^{64L}$ and agrees with threshold
        $\tfrac12$.
  \item \textbf{Involution sandwich.}
        Overlapping pairs are rejected; applying a valid pair-set twice
        is the identity on the alphabet.
  \item \textbf{Freeze.}
        Coordinates in $M$ do not contribute to $\pi$.
\end{enumerate}
\end{theorem}

\paragraph{Proof (machine-checked).}
Each clause is a \texttt{TensorLUTFormal.check*} in
\texttt{TensorLUTFormal.swift}, aggregated by
\texttt{TensorLUTFormal.certificate()}.

\reproduced{C25}{Theorem~1 corollary: emitter--discrete agreement
($\pi=0\Rightarrow E(w)$ recovers INIT bits) and involution under freeze
(\texttt{mutatedPreserving} retains frozen pairs).
\texttt{TensorLUTFormal.corollaryCertificate()}.
Reproduce:
\texttt{swift test -c release --filter testTensorLUTFormalCorollaryCertificate}.}

\begin{theorem}[TensorLUT emitter / freeze corollary; \corpusthm{1} corollary]
\label{thm:tensorlut-corollary}
Under Theorem~\ref{thm:tensorlut} hypotheses plus emit-via-$E$ and
\texttt{mutatedPreserving}:
if $\pi(w)=0$ then $E(w)$ recovers the binary INIT;
every freeze-preserving genotype remains a partial involution containing
the frozen pairs.
\end{theorem}

Still not melt completeness for arbitrary netlists.
The named lemmas of Theorem~\ref{thm:tensorlut} are
\texttt{discretenessPenalty},
\texttt{cryptoFitnessMSE},
\texttt{combinedObjective},
\texttt{emitterThreshold},
\texttt{involutionSandwich},
\texttt{freezeMask}.
Source path: \corpus{Sources/HELUTCore/TensorLUTFormal.swift}.

\begin{warning}
Theorem~1 does \emph{not} prove recovery of arbitrary keys; a U-534 /
P1030680 plaintext; or that melt is complete for all netlists.
Shatter / hold grades remain empirical (\cid{8}, \cid{9}, \hid{6}).
\end{warning}

\reproduced{C44}{Theorem~1$''$ melt--freeze--snap on a coordinate-separable
Boolean interpolant: unique maximizer $w=t$, snap basin
$\lvert w_i-t_i\rvert<1/2$, wrong freeze blocks $F=0$.
\texttt{TensorLUTFormal.meltFreezeSnapCertificate()}.
Reproduce:
\texttt{swift test -c release --filter testTensorLUTMeltFreezeSnapCertificate}.
Not multi-LUT topological melt.}

\reproduced{C48}{Two-LUT cascade $y=(a\land b)\oplus c$: melt from $0.5$,
snap, emit two LUT6 Verilog cells. Snapped INIT matches all 8 Boolean
corners (a dual interpolant such as NAND$+$XNOR is allowed).
\texttt{testTwoLUTCascadeMeltFreezeSnapEmit}. Not Yosys re-synth.}

\begin{theorem}[TensorLUT melt--freeze--snap, separable interpolant; \corpusthm{1$''$}]
\label{thm:tensorlut-snap}
Under a fully observed 1-LUT (each used INIT address an independent Boolean
target $t$), $F(w)=-\lVert w-t\rVert_2^2-\lambda\pi(w)$ has unique maximizer
$w=t$; the emitter recovers $t$ on the open cube
$\lvert w_i-t_i\rvert<1/2$; a freeze away from $t$ makes $F=0$ unreachable.
\end{theorem}

\section{Five-cell test for Theorem 1}

\begin{center}
\small
\begin{tabular}{@{}lp{0.68\textwidth}@{}}
\toprule
Cell & Artifact \\
\midrule
Proof & Theorem~\ref{thm:tensorlut} $+$ \texttt{TensorLUTFormalCertificate} \\
Table & Lemmas vs \texttt{holds} on the certificate; campaign grades in
        \corpus{BREAK\_P1030680.md} §21 \\
Metric & $F_{\mathrm{crypto}}=0$ on unmutated baseline; $\pi=0$ on binary INIT \\
Examples & M4 baseline emit; freeze-core involution; blind 3-pair PASS \\
Application & Stecker search that cannot propose a non-reciprocal map \\
\bottomrule
\end{tabular}
\end{center}

\begin{exercise}[Read a certificate]
Run the reproduce filter (Appendix~\ref{app:repro}) or, on the theory track,
read \texttt{TensorLUTFormal.certificate()} and list the six
\texttt{holds} bits.
For each lemma, write one sentence a journalist may print and one they
may not.
\end{exercise}

\section{Shatter versus hold (empirical)}

A full cold-start on a large sequential cipher discovers continuous
shortcuts that \emph{shatter} when $\lambda$ rises: $F_{\mathrm{crypto}}$
looks perfect in float and dies on the Boolean cube.
Targeted melting freezes known-good LUT blocks (\texttt{freezeMask}) and
evolves only the unknown cone---clause~6 of Theorem~\ref{thm:tensorlut}
is why frozen coordinates do not pay $\pi$.
Early resignation reboots doomed lineages rather than squeezing a
reciprocal shortcut to the end.

\begin{definition}[Grade]
Under a published $\lambda$ schedule, a lineage \textbf{holds} if the
emitted binary netlist preserves $F_{\mathrm{crypto}}=0$ (or a stated
bar). It \textbf{shatters} if discreteness kills the continuous solution.
Both grades are science. Shatter is not a failed demo; it is evidence
about shortcuts. It is \emph{not} a clause of Theorem~\ref{thm:tensorlut}.
\end{definition}

\reproduced{C8, C9}{TensorLUT M4 baseline $F_{\mathrm{crypto}}=0$ plus Verilog
emit; Stecker involution sandwich with blind 3-pair PASS.
Method claims, not a U-534 / P1030680 plaintext.}

\section{Involution sandwich}

Stecker genotypes are partial involutions (disjoint pairs) by construction
(Theorem~\ref{thm:tensorlut}, clause~5).
The GA cannot propose a non-reciprocal map.
The 16 cone LUTs around an identity plugboard cannot invent letter swaps by
INIT melt; the sandwich injects $S(\mathrm{CT})$ into a frozen core and
scores soft PT against bits of $S(P)$.
Reciprocity is structural.

\begin{warning}
Short cribs leave unused pairs unconstrained and can admit false $F=0$
steckers. Grade active-map agreement on $\mathrm{CT}\cup\mathrm{PT}$, use
parsimony, grow the pair budget on plateau, and allow soft freeze $+$ thaw.
TensorLUT grades are not a campaign decrypt (\claimkind{N}, \hid{6}).
\end{warning}

\section{Adversarial synthesis as a compiler loop}

Forward: silicon $\to$ uniform float tensors.
Reverse: cooled tensors $\to$ INIT hex $\to$ Verilog LUT instantiations.
The loop is a generative compiler in \emph{this} laboratory: cooled tensors
emit LUT6 Verilog, not a bitstream onto a physical FPGA.
Theorem~\ref{thm:tensorlut} is the invariant the loop is forced to obey.
It is not a guarantee that the loop finds a key.

\begin{exercise}
Why must $\pi(w)=\sum w_i(1-w_i)$ vanish on binary INITs, and why is a
fitness that ignores $\pi$ allowed to ``solve'' a cipher that no FPGA can
implement?
Give a two-LUT numerical cartoon.
Which clause of Theorem~\ref{thm:tensorlut} did you just unpack?
\end{exercise}

\begin{exercise}[Seminar]
Theorem~\ref{thm:tensorlut} is structural, not complete.
Write the statement of a completeness theorem you wish were true (melt
recovers $X$ on netlist class $Y$ under schedule $Z$), and the
counterexample shape you would try first.
Do not upgrade shatter/hold grades into that theorem.
\end{exercise}

\nonclaimbox{Pillar~II does not break P1030680, does not replace Welchman
search, and does not make TensorLUT a PBS. Parallel research; campaign
fitness stays cleartext.
\cid{19} is a theorem about the objective and the involution contract,
not about plaintext.}
